\noindent
\textbf{Egregi Professori Vardanega e Cardin,}
\vspace{0.3cm}
\noindent
Con la presente il gruppo \textbf{Git Happens} desidera presentare la propria candidatura per la realizzazione del progetto relativo al capitolato \textbf{C8 -- Sonar: Documentazione vivente di uno store Shopify}, sottoposto dall'azienda \textbf{Halue S.r.l. Societ\`a Benefit}.
\vspace{0.2cm}
\noindent
La decisione di concorrere per tale capitolato \`e maturata a seguito di una analisi comparativa tra i progetti rimasti disponibili e degli approfondimenti condotti con i rispettivi proponenti.
\vspace{0.4cm}
\section*{Documentazione di Riferimento}
La documentazione che ha contribuito a tale decisione \`e costituita da:
\begin{enumerate}
    \item \textbf{Verbali Interni:}
    \begin{itemize}
        \item Verbale Interno 2026/10/16
        \item Verbale Interno 2026/10/20
        \item Verbale Interno 2026/10/30
        \item Verbale Interno 2026/11/03
    \end{itemize}
    \item \textbf{Verbali Esterni:}
    \begin{itemize}
        \item Verbale Esterno 2026/10/20
        \item Verbale Esterno 2026/10/22
        \item Verbale Esterno 2026/10/26
        \item Verbale Esterno 2026/10/27
    \end{itemize}
    \item \textbf{Analisi dei Capitolati e Preventivo Costi}
\end{enumerate}
\noindent
Per maggiori informazioni si rimanda ai documenti visionabili sul \href{https://githappens-team.github.io/}{sito web ufficiale del gruppo}.
\vspace{0.4cm}
\section*{Motivazioni della Scelta}
A valle dell'analisi dei capitolati e del positivo riscontro ottenuto nel colloquio del 26 ottobre 2026, il team ha individuato in Sonar la proposta pi\`u stimolante e ad alto valore formativo:
\begin{itemize}
    \item \textbf{Chiarezza dei requisiti e tracciabilit\`a:} Il capitolato si distingue per la massima chiarezza e verificabilit\`a: requisiti con identificativi univoci, classi di priorit\`a ben definite e criteri di accettazione numerici. Questo consente di impostare fin da subito un tracciamento rigoroso tra requisiti, test e codice, in piena aderenza agli obiettivi didattici del corso.
    
    \item \textbf{Architettura a pipeline e modularit\`a:} L'architettura a pipeline descritta nel capitolato e approfondita nel colloquio tecnico (acquisizione tramite \textit{Shopify Admin GraphQL API}, normalizzazione in modello a grafo, analisi e presentazione) assegna a ogni livello responsabilit\`a e interfacce chiare. Questa struttura favorisce la separazione delle responsabilit\`a (\textit{SoC}), lo sviluppo in parallelo tra i sette membri e test d'integrazione ben isolati.
    
    \item \textbf{Stack tecnologico moderno e riproducibile:} Il progetto \`e interamente software e si sviluppa su un ecosistema moderno e altamente spendibile: \textit{TypeScript}, \textit{Next.js}, \textit{Supabase (PostgreSQL)} e le funzionalit\`a serverless di \textit{Vercel (Workflows, Blob, AI SDK)}. L'intero ambiente di lavoro \`e accessibile su piani gratuiti ed \`e riproducibile da ciascun componente del gruppo.
    
    \item \textbf{Trasparenza e supporto del proponente (Halue S.r.l.):} Il colloquio del 26 ottobre 2026 con il referente Valentino Baraldo ha confermato condizioni ideali per la collaborazione: grande apertura tecnica, piena autonomia nelle decisioni progettuali, disponibilit\`a a rinegoziare la portata dei requisiti in caso di necessit\`a e l'offerta di sessioni di consultazione con sviluppatori esperti di Shopify interni all'azienda.
\end{itemize}
\vspace{0.2cm}
\noindent
Per le motivazioni dettagliate si rimanda al documento \textit{Valutazione Capitolati} e al \href{https://githappens-team.github.io/content/verbali/esterni/2026-10-26/2026-10-26.pdf}{Verbale Esterno 2026/10/26}.
\vspace{0.4cm}
\section*{Costi e Tempistiche}
Come dettagliato nel documento \textit{Piano di Progetto e Preventivo Costi}, il gruppo stima l'impegno economico e la data di consegna finale come segue:
\begin{itemize}
    \item \textbf{Scadenza prevista per la consegna:} 26 Marzo 2027
    \item \textbf{Preventivo costi complessivo:} \geneuro\,12.845,00 (93 ore produttive per ciascun componente)
\end{itemize}
\vspace{0.6cm}
\section*{Composizione del Gruppo}
Di seguito i membri del gruppo \textbf{Git Happens}:
\vspace{0.3cm}

\begin{tabularx}{\textwidth}{@{} X l @{}}
    \toprule
    \textbf{Nominativo} & \textbf{Matricola} \\
    \midrule
    Diana Grosselle (Responsabile) & 2137985 \\
    \addlinespace[0.15cm]
    Cristiano Antoniazzo & 2138489 \\
    \addlinespace[0.15cm]
    Gabriele Gambarro & 2147979 \\
    \addlinespace[0.15cm]
    Gianmaria Grinovero & 2138007 \\
    \addlinespace[0.15cm]
    Tommaso Caon & 2111933 \\
    \addlinespace[0.15cm]
    Antonio Hai Xu Tang & 2111017 \\
    \addlinespace[0.15cm]
    Filippo Lissandrin & 2137980 \\
    \bottomrule
\end{tabularx}
\vspace{0.3cm}


Il team \textbf{Git Happens} ringrazia per l'attenzione.
\vspace{0.4cm}
\noindent
Cordiali saluti, \\
\textbf{Git Happens}
\vspace{0.5cm}
\noindent
\textit{Padova, 3 Novembre 2026}
